\section{Set Proofs: Element Chasing in Detail}
\chapterline

In Chapter~10, we introduced sets, membership ($\in$), and operations such as union, intersection, and set difference. When encountering an identity like
\[
A \cap (B \cup C) = (A \cap B) \cup (A \cap C),
\]
it is tempting to sketch overlapping circles and declare the result obvious.

Venn diagrams build valuable visual intuition, but in pure mathematics, a picture cannot serve as the final proof. For two or three sets, circles are convenient. But what happens when you have four sets? As John Venn discovered, it is geometrically impossible to draw four mutually intersecting sets with circles---you must switch to ellipses or interlocking curves. For five sets, diagrams become unwieldy. And for infinite families of sets in analysis, sketches fail completely. Visual diagrams can also conceal edge cases involving the empty set $\emptyset$.

To verify set identities rigorously, mathematicians use \emph{element chasing}. The idea is straightforward: take an arbitrary element $x$ from one set and follow its logical properties until showing that $x$ belongs to the second set.

\subsection{The Architecture of Set Containment}

A set is defined entirely by which elements belong to it, not by the names or formulas we use to describe it---a principle known as \emph{extensionality}. If two sets $A$ and $B$ contain the exact same items, they are identical. 

To prove that one set is contained in another, we recall the definition of subset inclusion:
\[
A \subseteq B \iff (\forall x)(x \in A \implies x \in B).
\]
Proving $A \subseteq B$ is not about manipulating entire collections all at once. It is a conditional proof about an arbitrary element:
\begin{enumerate}[label=\arabic*.]
\item Begin by introducing an arbitrary element: ``Let $x \in A$.''
\item Unpack what $x \in A$ means using definitions, hypotheses, and logic.
\item Conclude that $x$ must satisfy the membership condition for $B$: ``Therefore $x \in B$.''
\end{enumerate}

To prove that two sets $A$ and $B$ are \emph{equal}, we establish mutual containment in both directions. Two sets are identical if and only if neither set contains an element foreign to the other:
\[
A = B \iff (A \subseteq B \text{ and } B \subseteq A).
\]

\begin{figure}[htbp]
\centering
\begin{tikzpicture}[scale=0.92, >=Latex]
  % Panel 1: Double containment architecture
  \begin{scope}[xshift=0cm]
    \draw[rounded corners=4pt, fill=gray!3, draw=gray!25] (-0.4, -0.7) rectangle (5.4, 3.5);
    \node[anchor=south, font=\small\bfseries] at (2.5, 3.6) {Two-Way Containment: $A = B$};
    
    % Sets A and B
    \draw[thick, sourceblue, fill=sourceblue!12] (1.6, 1.4) ellipse (1.3cm and 1.1cm);
    \draw[thick, targetgreen, fill=targetgreen!12] (3.4, 1.4) ellipse (1.3cm and 1.1cm);
    
    \node[sourceblue!90!black, font=\bfseries] at (0.8, 1.4) {$A$};
    \node[targetgreen!90!black, font=\bfseries] at (4.2, 1.4) {$B$};
    
    % Direction arrows
    \draw[->, very thick, sourceblue!85!black] (1.7, 2.65) to[bend left=22] node[above=2pt, font=\footnotesize, text=sourceblue!90!black] {Direction 1: $x \in A \implies x \in B$} (3.3, 2.65);
    \draw[->, very thick, targetgreen!85!black] (3.3, 0.15) to[bend left=22] node[below=2pt, font=\footnotesize, text=targetgreen!90!black] {Direction 2: $y \in B \implies y \in A$} (1.7, 0.15);
  \end{scope}

  % Panel 2: Element chase through distribution
  \begin{scope}[xshift=6.4cm]
    \draw[rounded corners=4pt, fill=gray!3, draw=gray!25] (-0.4, -0.7) rectangle (5.8, 3.5);
    \node[anchor=south, font=\small\bfseries] at (2.7, 3.6) {Element Chasing: $x \in A \cap (B \cup C)$};

    % Tracking path
    \node[draw=sourceblue, fill=sourceblue!15, rounded corners=3pt, font=\footnotesize, inner sep=4pt] (start) at (2.7, 2.65) {$x \in A \text{ and } x \in (B \cup C)$};
    \node[draw=structurepurple, fill=structurepurple!15, rounded corners=3pt, font=\footnotesize, inner sep=4pt] (branch) at (2.7, 1.4) {$x \in A \text{ and } (x \in B \text{ or } x \in C)$};
    
    \node[draw=targetgreen, fill=targetgreen!15, rounded corners=3pt, font=\footnotesize, inner sep=4pt] (left) at (1.2, 0.15) {$x \in A \cap B$};
    \node[draw=quotientorange, fill=quotientorange!15, rounded corners=3pt, font=\footnotesize, inner sep=4pt] (right) at (4.2, 0.15) {$x \in A \cap C$};
    
    \draw[->, thick, sourceblue] (start) -- (branch);
    \draw[->, thick, targetgreen] (branch) -- node[left=2pt, font=\scriptsize] {case $x \in B$} (left);
    \draw[->, thick, quotientorange] (branch) -- node[right=2pt, font=\scriptsize] {case $x \in C$} (right);
  \end{scope}
\end{tikzpicture}
\caption{The two foundational architectures of set proofs. \textbf{Left}: Demonstrating set equality $A = B$ requires establishing mutual subset inclusion in two independent directions---first showing that every $x \in A$ must belong to $B$, and then showing that every $y \in B$ must belong to $A$. \textbf{Right}: Chasing an element through compound logical conditions. When the hypothesis contains a union ($B \cup C$), logic presents an ``or'' (disjunction $\lor$), causing the chase to branch into exhaustive cases that each reach the target set.}
\label{11:fig-double-containment}
\end{figure}

\begin{intuitionbox}
Think of element chasing as reasoning about an arbitrary element $x$. You do not pick a specific number like $x = 0$ or $x = 5$. You write ``Let $x \in A$,'' using only the properties guaranteed by membership in $A$. If you can deduce from those properties that $x \in B$, then every element of $A$ must belong to $B$.
\end{intuitionbox}

\subsection{A Warm-Up: Intersection Subsets}

Let us start with the most basic structural property of intersections.

\begin{theorem}
For any sets $A$ and $B$, $A \cap B \subseteq A$.
\end{theorem}

\begin{proof}
Let $x \in A \cap B$ be arbitrary. 

By definition of intersection, $x \in A \cap B$ means that $x \in A$ and $x \in B$. 

In formal logic, a statement joined by ``and'' is called a \emph{conjunction} ($P \land Q$). From the conjunction $x \in A \land x \in B$, the single statement $x \in A$ follows immediately by simplification (if both facts are true, each one is true on its own). 

Since $x$ was an arbitrary element of $A \cap B$, every element of $A \cap B$ belongs to $A$. Thus $A \cap B \subseteq A$.
\end{proof}

Notice how clean the flow is: we assumed $x \in A \cap B$, unwrapped the definition into a logical statement, extracted what was needed, and concluded containment.

\begin{intuitionbox}
\textbf{Pause and think:} Why couldn't we prove $A \cap B \subseteq A$ by testing concrete sets, say $A = \{1, 2, 3\}$ and $B = \{2, 3, 4\}$, calculating $A \cap B = \{2, 3\}$, and observing that both $2$ and $3$ belong to $A$? Because testing concrete numbers only checks that specific example. To establish a universal claim ($(\forall x)(x \in A \cap B \implies x \in A)$), the element $x$ must be generic, possessing only the properties guaranteed by membership in $A \cap B$.
\end{intuitionbox}

\subsection{Set Difference as Complementation}

Next, consider the relationship between set difference and the complement. Before we can talk about a complement, we must specify an ambient \emph{universal set} $U$---the overarching background universe of objects under discussion (such as the real numbers $\R$ or the integers $\Z$). 

The \emph{complement} of a set $B$, denoted $B^c$, consists of all objects in $U$ that lie strictly outside $B$:
\[
B^c = \{x \in U : x \notin B\}.
\]
The \emph{set difference} $A \setminus B$ (read ``$A$ minus $B$'') consists of all elements that belong to $A$ with anything belonging to $B$ discarded:
\[
A \setminus B = \{x \in U : x \in A \text{ and } x \notin B\}.
\]
How do these two concepts interact? Removing $B$ from $A$ turns out to be exactly the same as intersecting $A$ with everything outside $B$.

\begin{theorem}
Let $A$ and $B$ be subsets of a universal set $U$. Then
\[
A \setminus B = A \cap B^c.
\]
\end{theorem}

Figure~\ref{11:fig-difference-complementation} illustrates this identity: whether you start with $A$ and remove $B$, or take the intersection of $A$ with the exterior of $B$ ($B^c$), the resulting territory is identical.

\begin{figure}[htbp]
\centering
\begin{tikzpicture}[
    font=\small,
    univ/.style={draw=gray!40, fill=gray!2, rounded corners=5pt, thick}
]

% Panel 1: A \ B
\begin{scope}[xshift=0cm]
    \draw[univ] (-0.3, -0.6) rectangle (5.3, 3.6);

    % Fills first
    \fill[sourceblue!25] (1.8, 1.3) circle (1.05);
    \fill[white] (3.2, 1.3) circle (1.05);

    % Outlines
    \draw[thick, sourceblue!80] (1.8, 1.3) circle (1.05);
    \draw[thick, gray!50] (3.2, 1.3) circle (1.05);

    % Labels last
    \node[anchor=north, font=\footnotesize\bfseries, text=sourceblue!90!black] at (2.5, 3.4) {Set Difference: $A \setminus B$};
    \node[font=\footnotesize\bfseries, text=sourceblue!90!black] at (1.2, 1.3) {$A$};
    \node[font=\footnotesize\bfseries, text=gray!60!black] at (3.8, 1.3) {$B$};

    \node[anchor=south, font=\scriptsize, text=gray!80!black] at (2.5, -0.45) {Start with $A$, discard elements in $B$};
\end{scope}

% Equals symbol
\node[font=\LARGE\bfseries, text=gray!60!black] at (5.9, 1.3) {$=$};

% Panel 2: A \cap B^c
\begin{scope}[xshift=6.5cm]
    % Background box with B^c shading first
    \draw[draw=gray!40, fill=quotientorange!10, rounded corners=5pt, thick] (-0.3, -0.6) rectangle (5.3, 3.6);
    \fill[white] (3.2, 1.3) circle (1.05);

    % Intersection A \cap B^c fill
    \begin{scope}
        \clip (1.8, 1.3) circle (1.05);
        \fill[targetgreen!25] (1.8, 1.3) circle (1.05);
        \fill[white] (3.2, 1.3) circle (1.05);
    \end{scope}

    % Outlines
    \draw[thick, sourceblue!80] (1.8, 1.3) circle (1.05);
    \draw[thick, gray!50] (3.2, 1.3) circle (1.05);

    % Labels last
    \node[anchor=north, font=\footnotesize\bfseries, text=targetgreen!90!black] at (2.5, 3.4) {Complementation: $A \cap B^c$};
    \node[font=\footnotesize\bfseries, text=sourceblue!90!black] at (1.2, 1.3) {$A$};
    \node[font=\footnotesize\bfseries, text=gray!60!black] at (3.8, 1.3) {$B$};

    \node[anchor=south, font=\scriptsize, text=gray!80!black] at (2.5, -0.45) {Overlap of $A$ with exterior $B^c$ (shaded)};
\end{scope}

\end{tikzpicture}
\caption{Geometric equivalence of set difference and complementation ($A \setminus B = A \cap B^c$). On the left, we begin with $A$ and remove all points of $B$. On the right, we intersect $A$ with the exterior region $B^c$ (shaded in light orange). Both operations isolate the identical blue-green territory.}\label{11:fig-difference-complementation}
\end{figure}

\begin{proof}
We will prove this equality by demonstrating that an arbitrary element $x \in U$ belongs to the left-hand side if and only if it belongs to the right-hand side.

Let $x \in U$ be arbitrary. Then:
\begin{align*}
x \in A \setminus B 
&\iff x \in A \text{ and } x \notin B && \text{(definition of set difference)} \\
&\iff x \in A \text{ and } x \in B^c && \text{(definition of complement)} \\
&\iff x \in A \cap B^c && \text{(definition of intersection)}.
\end{align*}
Because each step is a bidirectional logical equivalence ($\iff$), the two sets contain exactly the same elements. Hence $A \setminus B = A \cap B^c$.
\end{proof}

\begin{trapbox}
Chaining if-and-only-if ($\iff$) statements is elegant when every step is truly reversible. But be careful: if even a single step only holds in one direction (for instance, $x \in S \implies x \in S \cup T$, which cannot be reversed), the whole two-way proof collapses. When in doubt, split the proof into two explicit halves: first assume $x$ is in the left set and deduce it is in the right, then assume $x$ is in the right set and deduce it is in the left.
\end{trapbox}

\subsection{Distributing Intersection Over Union}

In basic algebra, multiplication distributes across addition: $a \cdot (b + c) = (a \cdot b) + (a \cdot c)$. In set theory, intersection distributes across union in the same way:
\[
A \cap (B \cup C) = (A \cap B) \cup (A \cap C).
\]
Why does proving this identity require a case-by-case argument? When an element $x$ belongs to a union $B \cup C$, logic presents us with a \emph{disjunction} ($P \lor Q$)---the condition that $x \in B$ or $x \in C$. Unlike a conjunction (``and''), which allows us to hold both facts simultaneously, a disjunction offers two alternative possibilities. Since we do not know in advance which condition $x$ satisfies, we consider each case separately and verify that in both cases $x$ belongs to $(A \cap B) \cup (A \cap C)$.

\begin{theorem}
For all sets $A, B$, and $C$,
\[
A \cap (B \cup C) = (A \cap B) \cup (A \cap C).
\]
\end{theorem}

Figure~\ref{11:fig-distributive-branching} contrasts the spatial intuition with the formal logical architecture: because union introduces an ``or'' condition, the forward element chase must split into two parallel cases that independently reach the destination.

\begin{figure}[htbp]
\centering
\begin{tikzpicture}[
    font=\small,
    univ/.style={draw=gray!40, fill=gray!2, rounded corners=5pt, thick},
    boxnode/.style={draw=gray!35, fill=white, rounded corners=4pt, inner sep=4pt, align=center}
]

% Panel 1: Geometric Venn Diagram
\begin{scope}[xshift=0cm]
    \draw[univ] (-0.3, -0.7) rectangle (4.3, 4.3);

    % Fills first
    \begin{scope}
        \clip (2.0, 2.3) circle (1.05);
        \fill[sourceblue!25] (1.3, 1.2) circle (1.05);
        \fill[targetgreen!25] (2.7, 1.2) circle (1.05);
        % Overlap region
        \begin{scope}
            \clip (1.3, 1.2) circle (1.05);
            \fill[sourceblue!20!targetgreen!20] (2.7, 1.2) circle (1.05);
        \end{scope}
    \end{scope}

    % Outlines
    \draw[thick, sourceblue!80] (2.0, 2.3) circle (1.05);
    \draw[thick, gray!55] (1.3, 1.2) circle (1.05);
    \draw[thick, gray!55] (2.7, 1.2) circle (1.05);

    % Labels last
    \node[anchor=north, font=\footnotesize\bfseries, text=sourceblue!90!black] at (2.0, 4.05) {Venn Territory};
    \node[font=\footnotesize\bfseries, text=sourceblue!90!black] at (2.0, 2.9) {$A$};
    \node[font=\footnotesize\bfseries, text=gray!60!black] at (0.6, 1.0) {$B$};
    \node[font=\footnotesize\bfseries, text=gray!60!black] at (3.4, 1.0) {$C$};

    \node[font=\scriptsize\bfseries, text=sourceblue!90!black] at (1.4, 1.9) {$A \cap B$};
    \node[font=\scriptsize\bfseries, text=targetgreen!90!black] at (2.6, 1.9) {$A \cap C$};

    \node[anchor=south, font=\scriptsize, text=gray!80!black] at (2.0, -0.55) {Two shaded lobes form $A \cap (B \cup C)$};
\end{scope}

% Panel 2: Logical Branching Diagram
\begin{scope}[xshift=4.9cm]
    \draw[univ] (-0.2, -0.7) rectangle (8.4, 4.3);

    \node[anchor=north, font=\footnotesize\bfseries, text=structurepurple!90!black] at (4.1, 4.05) {Disjunction Forces a Case Split};

    % Top node
    \node[boxnode, draw=sourceblue!50] (top) at (4.1, 2.85) {
        \textbf{Hypothesis:} $x \in A \cap (B \cup C)$ \\[2pt]
        {\scriptsize $x \in A$ \textbf{and} ($x \in B$ or $x \in C$)}
    };

    % Case 1 and Case 2 nodes
    \node[boxnode, draw=sourceblue!40] (c1) at (1.9, 1.45) {
        \textbf{Case 1: $x \in B$} \\[2pt]
        $x \in A$ and $x \in B$ \\[1pt]
        $\implies x \in A \cap B$
    };

    \node[boxnode, draw=targetgreen!50] (c2) at (6.3, 1.45) {
        \textbf{Case 2: $x \in C$} \\[2pt]
        $x \in A$ and $x \in C$ \\[1pt]
        $\implies x \in A \cap C$
    };

    % Conclusion node
    \node[boxnode, draw=targetgreen!60, fill=targetgreen!8] (bot) at (4.1, 0.05) {
        \textbf{Union Property:} in either case, \\[1pt]
        $x \in (A \cap B) \cup (A \cap C)$
    };

    % Arrows between nodes
    \draw[->, thick, sourceblue!70] (top.south -| c1.north) -- (c1.north);
    \draw[->, thick, targetgreen!70] (top.south -| c2.north) -- (c2.north);
    \draw[->, thick, sourceblue!70] (c1.south) -- (2.6, 0.45);
    \draw[->, thick, targetgreen!70] (c2.south) -- (5.6, 0.45);
\end{scope}

\end{tikzpicture}
\caption{Distributive law for sets ($A \cap (B \cup C) = (A \cap B) \cup (A \cap C)$). Left: the geometric overlap consists of two distinct regions, $A \cap B$ and $A \cap C$. Right: the element chase structure. Because membership in a union presents an ``or'' condition, the argument must branch into two exhaustive cases, each independently certifying membership in the target union.}\label{11:fig-distributive-branching}
\end{figure}

\begin{proof}
We prove mutual containment: first $A \cap (B \cup C) \subseteq (A \cap B) \cup (A \cap C)$, and then the reverse inclusion.

\textbf{Direction 1: $(\subseteq)$.} Let $x \in A \cap (B \cup C)$ be arbitrary. 
By definition of intersection, $x \in A$ and $x \in B \cup C$. 
Unpacking the union, $x \in B \cup C$ means $x \in B$ or $x \in C$. We now argue by cases:
\begin{itemize}
\item \emph{Case 1: $x \in B$.} Since we already know $x \in A$, we have $x \in A$ and $x \in B$, which means $x \in A \cap B$. It follows that $x \in (A \cap B) \cup (A \cap C)$.
\item \emph{Case 2: $x \in C$.} Since $x \in A$, we have $x \in A$ and $x \in C$, so $x \in A \cap C$. Therefore $x \in (A \cap B) \cup (A \cap C)$.
\end{itemize}
In either case, $x \in (A \cap B) \cup (A \cap C)$. This establishes that
\[
A \cap (B \cup C) \subseteq (A \cap B) \cup (A \cap C).
\]

\textbf{Direction 2: $(\supseteq)$.} Now let $y \in (A \cap B) \cup (A \cap C)$ be arbitrary.
By definition of union, $y \in A \cap B$ or $y \in A \cap C$. Again, we examine the two possibilities:
\begin{itemize}
\item \emph{Case 1: $y \in A \cap B$.} Then $y \in A$ and $y \in B$. Since $y \in B$, we certainly have $y \in B \cup C$. Combining $y \in A$ with $y \in B \cup C$ gives $y \in A \cap (B \cup C)$.
\item \emph{Case 2: $y \in A \cap C$.} Then $y \in A$ and $y \in C$. Since $y \in C$, we have $y \in B \cup C$. Hence $y \in A \cap (B \cup C)$.
\end{itemize}
In both cases, $y \in A \cap (B \cup C)$. Thus
\[
(A \cap B) \cup (A \cap C) \subseteq A \cap (B \cup C).
\]

Having established containment in both directions, we conclude that $A \cap (B \cup C) = (A \cap B) \cup (A \cap C)$.
\end{proof}

\begin{historybox}
In the 1760s, Leonhard Euler introduced closed circular curves to illustrate categorical syllogisms in logic. A century later, in an 1880 paper titled ``On the Diagrammatic and Mechanical Representation of Propositions and Reasonings,'' the English logician John Venn refined Euler's idea by requiring that all possible combinations of inclusion and exclusion among the sets be visually represented, even building mechanical wooden devices with sliding gears to model logical deductions. For two or three sets, Venn's diagrams are wonderfully clear. 

Yet mathematics quickly ran into hard geometric and conceptual barriers. First, Venn showed that while three sets can be drawn with congruent circles, four mutually intersecting sets cannot: congruent circles can only form $14$ of the required $16$ regions, forcing one to use distorted ellipses or interlocking ribbons. For five or more sets, constructing symmetric Venn diagrams made of congruent curves is mathematically impossible unless the number of sets $n$ is prime (proven by Anthony Edwards and Branko Gr\"unbaum).

Second, visual diagrams break down completely when confronted with infinite collections of sets, such as the shrinking sequence of intervals $\bigcap_{n=1}^\infty (-1/n, 1/n) = \{0\}$ in real analysis, where there are infinitely many boundaries. 

Element chasing replaced spatial inspection with formal symbolic deduction: reasoning about an arbitrary element $x$ applies equally well to two sets, twenty sets, or an arbitrary family of sets.
\end{historybox}

\subsection{Power Sets and Inclusion: Chasing Across Layers}

Element chasing also works when the elements themselves are sets. The \emph{power set} of a set $A$, denoted $\pow(A)$, is the collection whose elements are all possible subsets of $A$:
\[
X \in \pow(A) \iff X \subseteq A.
\]
For example, if $A = \{1, 2\}$, its subsets are $\emptyset, \{1\}, \{2\}$, and $\{1, 2\}$, so $\pow(A) = \{\emptyset, \{1\}, \{2\}, \{1, 2\}\}$. 

When we prove an inclusion between power sets, the element being chased is itself a set.

\begin{theorem}
If $A \subseteq B$, then $\pow(A) \subseteq \pow(B)$.
\end{theorem}

\begin{proof}
Suppose $A \subseteq B$. We want to show $\pow(A) \subseteq \pow(B)$.

Let $X \in \pow(A)$ be arbitrary. By definition of the power set, this means $X \subseteq A$.

Now we need to show that $X \in \pow(B)$, which means showing $X \subseteq B$. 

To show $X \subseteq B$, let $z \in X$ be an arbitrary element. Since $X \subseteq A$, we have $z \in A$. Since $A \subseteq B$, we have $z \in B$. Since every $z \in X$ belongs to $B$, we have $X \subseteq B$.

Therefore $X \in \pow(B)$. Because $X \in \pow(A)$ was arbitrary, $\pow(A) \subseteq \pow(B)$.
\end{proof}

\begin{figure}[htbp]
\centering
\begin{tikzpicture}[scale=0.92, >=Latex]
  % Panel 1: Outer Level (Power Sets)
  \begin{scope}[xshift=0cm]
    \draw[rounded corners=4pt, fill=gray!3, draw=gray!25] (-0.4, -0.7) rectangle (5.6, 3.5);
    \node[anchor=south, font=\small\bfseries] at (2.6, 3.6) {Outer Level: Power Sets ($\in$)};

    % Big ellipse: P(B)
    \draw[thick, targetgreen, fill=targetgreen!8] (2.6, 1.4) ellipse (2.4cm and 1.4cm);
    \node[targetgreen!90!black, font=\bfseries\small] at (4.2, 2.05) {$\pow(B)$};

    % Smaller ellipse: P(A)
    \draw[thick, sourceblue, fill=sourceblue!15] (1.8, 1.3) ellipse (1.4cm and 0.95cm);
    \node[sourceblue!90!black, font=\bfseries\small] at (1.1, 1.8) {$\pow(A)$};

    % The element X inside P(A)
    \node[draw=structurepurple, fill=structurepurple!20, rounded corners=3pt, inner sep=3pt, font=\footnotesize\bfseries] (Xnode) at (2.2, 1.2) {Set $X$};
    \node[font=\scriptsize, structurepurple!90!black] at (2.2, 0.65) {$X \in \pow(A)$};
  \end{scope}

  % Zoom connector between panels
  \draw[->, very thick, dashed, structurepurple!80!black] (5.7, 1.6) -- node[above=2pt, font=\scriptsize\bfseries, text=structurepurple!90!black] {Zoom} (7.1, 1.6);

  % Panel 2: Inner Level (Points inside Sets)
  \begin{scope}[xshift=7.4cm]
    \draw[rounded corners=4pt, fill=gray!3, draw=gray!25] (-0.4, -0.7) rectangle (5.6, 3.5);
    \node[anchor=south, font=\small\bfseries] at (2.6, 3.6) {Inner Level: Point Chasing ($\subseteq$)};

    % Set B
    \draw[thick, targetgreen, fill=targetgreen!8] (2.6, 1.4) ellipse (2.3cm and 1.4cm);
    \node[targetgreen!90!black, font=\bfseries\small] at (4.2, 2.05) {$B$};

    % Set A
    \draw[thick, sourceblue, fill=sourceblue!15] (2.0, 1.3) ellipse (1.5cm and 1.05cm);
    \node[sourceblue!90!black, font=\bfseries\small] at (1.2, 1.8) {$A$};

    % Set X inside A
    \draw[thick, structurepurple, fill=structurepurple!25] (2.2, 1.2) ellipse (1.0cm and 0.65cm);
    \node[structurepurple!90!black, font=\scriptsize\bfseries] at (1.6, 1.4) {$X$};

    % Point z inside X
    \node[circle, fill=black, inner sep=1.4pt] (zpt) at (2.2, 1.1) {};
    \node[right=2pt, font=\scriptsize, text=black] at (2.2, 1.1) {$z \in X$};
    
    \node[font=\scriptsize, align=center] at (2.6, -0.3) {$z \in X \implies z \in A \implies z \in B$\\$\therefore X \subseteq B \iff X \in \pow(B)$};
  \end{scope}
\end{tikzpicture}
\caption{Two levels of element chasing in power sets. \textbf{Left}: At the outer level, the element is a set $X$ satisfying $X \in \pow(A)$. \textbf{Right}: To show $X \subseteq B$, we consider an arbitrary element $z \in X$. Since $X \subseteq A$ and $A \subseteq B$, we deduce $z \in B$, establishing $X \subseteq B$ and therefore $X \in \pow(B)$.}
\label{11:fig-power-set-levels}
\end{figure}

\begin{labbox}
Notice the two layers of element chasing in the proof above (illustrated in Figure~\ref{11:fig-power-set-levels})! At the outer level, the element is a set $X$ moving from $\pow(A)$ to $\pow(B)$. At the inner level, the element is a point $z$ moving from $X$ to $B$. Keeping these levels distinct prevents the common beginner mistake of confusing $\in$ with $\subseteq$.
\end{labbox}

\begin{intuitionbox}
\textbf{Pause and think:} What happens in the extreme case where $A = \emptyset$? Does Theorem~4 still hold? Yes! The empty set has only one subset: itself ($\pow(\emptyset) = \{\emptyset\}$). If $A = \emptyset$, then $A \subseteq B$ is vacuously true for every set $B$, and since $\emptyset \subseteq B$, we have $\emptyset \in \pow(B)$. Thus $\{\emptyset\} \subseteq \pow(B)$, confirming that $\pow(\emptyset) \subseteq \pow(B)$ as required.
\end{intuitionbox}

\subsection{Scope of element chasing and the transition to functions}

Element chasing provides a direct method to verify set identities, Boolean laws, and subset inclusions. It replaces visual diagrams with rigorous logical deduction.

Sets and element chasing describe static collections and containment, but most of mathematics concerns how elements of one set are sent to elements of another---and whether those transformations preserve algebraic, geometric, or topological structure. In Chapter~12 (\emph{Functions: Structured Promises, Not Just Formulas}), we study functions, the precise rules that map inputs in one set to outputs in another.

\subsection{Exercises}

\subsubsection*{Warm-Up Mental Checks}
\begin{enumerate}[label=\arabic*.]
\item \textbf{The One-Way Trap:} A student writes: ``$x \in A \cup B \iff x \in A$, therefore $A \cup B \subseteq A$.'' Where is the exact logical flaw? Can you give a simple counterexample with numbers where $x \in A \cup B$ is true but $x \in A$ is false?
\item \textbf{Power Set Distinctions:} Let $A = \{1, 2\}$. For each of the following statements, decide whether it is true or false, and explain the difference between them in your own words:
\begin{enumerate}[label=(\alph*)]
\item $\emptyset \in \pow(A)$
\item $\emptyset \subseteq \pow(A)$
\item $\{\emptyset\} \subseteq \pow(A)$
\end{enumerate}
\item \textbf{Proving Emptiness:} Suppose you want to prove that a set $S$ is empty ($S = \emptyset$). How can you adapt the element-chasing method to achieve this? (Hint: Assume for contradiction that there exists some $x \in S$, and chase this hypothetical element until you deduce a logical impossibility.)
\end{enumerate}

\subsubsection*{Core Practice: Element-Chasing Proofs}
Write out complete proofs using full grammatical sentences, explicitly introducing every arbitrary element.
\begin{enumerate}[resume*]
\item Prove that for any sets $A$ and $B$, $A \subseteq A \cup B$.
\item Prove that $A \cap (B \setminus C) = (A \cap B) \setminus C$.
\item Prove De Morgan's Law for set difference across unions:
\[
A \setminus (B \cup C) = (A \setminus B) \cap (A \setminus C).
\]
\item Prove the companion De Morgan identity for set difference across intersections:
\[
A \setminus (B \cap C) = (A \setminus B) \cup (A \setminus C).
\]
\end{enumerate}

\subsubsection*{Deeper Explorations}
\begin{enumerate}[resume*]
\item \textbf{Power Sets and Intersections:} Prove that for all sets $A$ and $B$,
\[
\pow(A \cap B) = \pow(A) \cap \pow(B).
\]
\item \textbf{Power Sets and Unions:} What happens if we replace intersection with union in the previous problem?
\begin{enumerate}[label=(\alph*)]
\item Prove that $\pow(A) \cup \pow(B) \subseteq \pow(A \cup B)$ always holds.
\item Give an explicit counterexample with small finite sets to prove that the reverse containment, $\pow(A \cup B) \subseteq \pow(A) \cup \pow(B)$, is generally false.
\end{enumerate}
\item \textbf{Infinite Families:} Let $\{A_i\}_{i \in I}$ be an indexed family of sets (where the index set $I$ may be finite or infinite), and let $B$ be a set. Recall that an element belongs to an arbitrary intersection $\bigcap_{i \in I} S_i$ if and only if it belongs to $S_i$ for every $i \in I$. Prove that:
\[
\left( \bigcap_{i \in I} A_i \right) \cap B = \bigcap_{i \in I} (A_i \cap B).
\]
Notice how tracking an arbitrary element $x$ works for an arbitrary indexed family just as it does for two sets.
\end{enumerate}

\newpage
